\chapter{AI tạo sinh cho nghiên cứu}

\subsection{Tóm tắt}

Chương 2 khám phá cách AI tạo sinh đang định hình lại vòng đời nghiên cứu học thuật, từ hình thành ý tưởng và khám phá tài liệu đến xây dựng giả thuyết, lập kế hoạch phương pháp luận và thu thập dữ liệu. Bằng cách tăng cường các quy trình ở giai đoạn đầu thông qua các công cụ như GPT, Claude và Gemini, nhà nghiên cứu có thể hợp lý hóa việc phát triển khái niệm, phát hiện các kết nối liên ngành và tổng hợp tài liệu một cách nhanh chóng. AI tạo sinh hỗ trợ việc đề xuất giả thuyết và lựa chọn phương pháp, cung cấp hỗ trợ kỹ thuật cho lập trình và thu thập dữ liệu, đồng thời thúc đẩy sự tích hợp giữa các công cụ. Tuy nhiên, giá trị của nó không chỉ nằm ở tự động hóa mà ở sự cộng tác: khi được sử dụng một cách có tư duy phản biện, các hệ thống này nâng cao hiểu biết học thuật mà không làm tổn hại đến tính chặt chẽ của nghiên cứu. Chương này nhấn mạnh tầm quan trọng của việc kiểm chứng, tính minh bạch và sự thận trọng về mặt đạo đức trong thiết kế nghiên cứu có sự hỗ trợ của AI.

\subsection{Từ khóa}

AI tạo sinh, nghiên cứu học thuật, tổng quan tài liệu, hình thành giả thuyết, phương pháp luận nghiên cứu

Vòng đời nghiên cứu học thuật truyền thống diễn ra qua một chuỗi các giai đoạn liên kết với nhau, hình thành ý tưởng ban đầu, tổng quan tài liệu, hình thành giả thuyết, thu thập dữ liệu, phân tích, viết và phổ biến kết quả (Creswell, 2014). Mỗi bước đều đòi hỏi nhà nghiên cứu phải dành nhiều nỗ lực trí tuệ, tư duy sáng tạo, sự chính xác về phương pháp luận và tính chặt chẽ trong phân tích. Sự xuất hiện của AI tạo sinh (generative AI), đặc biệt là các mô hình ngôn ngữ lớn (LLM) tiên tiến như GPT (OpenAI), Claude (Anthropic), Llama (Meta) hay Gemini (Google), đang làm thay đổi căn bản cách các học giả tiếp cận những giai đoạn đầu của hoạt động nghiên cứu, mang lại những lối đi mới cho việc hình thành ý tưởng, xây dựng khung khái niệm và khám phá phương pháp luận.

Quy trình nghiên cứu truyền thống, tuy vững chắc về phương pháp luận, vẫn đối mặt với những hạn chế đáng kể trong thế giới học thuật biến đổi nhanh chóng ngày nay. Các nhà nghiên cứu phải xử lý một mạng lưới liên kết liên ngành ngày càng phức tạp, quản lý khối lượng tài liệu tăng theo cấp số nhân, và xác định những đóng góp mới trong các lĩnh vực ngày càng chuyên sâu. Các nghiên cứu gần đây cho thấy những nhà nghiên cứu sử dụng công cụ AI tạo sinh trong giai đoạn đầu của quá trình nghiên cứu báo cáo giảm tới 40\% thời gian dành cho tổng quan tài liệu ban đầu và phát triển khái niệm (Noy and Zhang, 2023). Những lợi ích về hiệu quả này đặc biệt rõ rệt trong nghiên cứu liên ngành, nơi các công cụ AI nổi trội trong việc nhận diện các mối liên hệ xuyên lĩnh vực mà nếu không thì có thể vẫn chưa được khám phá.

Các thảo luận học thuật gần đây nhấn mạnh rằng AI tạo sinh không đóng vai trò thay thế chuyên môn của nhà nghiên cứu, mà là một lực lượng bổ trợ giúp tăng cường năng lực trí tuệ của con người (van Dis et al., 2023; Rossi et al., 2024). Mối quan hệ cộng hưởng này đã được ghi nhận ở nhiều ngành, từ khoa học xã hội đến nghiên cứu y học, nơi việc thiết kế nghiên cứu có sự hỗ trợ của AI đã dẫn đến những câu hỏi nghiên cứu toàn diện và tinh tế hơn (Borger et al., 2023; Oksanen, 2024). Chẳng hạn, các nhà nghiên cứu có thể tận dụng công cụ AI tạo sinh để đưa ra nhiều góc nhìn về các vấn đề nghiên cứu, nhanh chóng tổng hợp tài liệu hiện có và xác định các hướng tiếp cận phương pháp luận tiềm năng, đồng thời vẫn duy trì sự giám sát có tính phản biện của con người đối với quá trình này.

Việc tích hợp AI tạo sinh phù hợp với thực tế của môi trường học thuật ngày nay: giàu dữ liệu, năng động và ngày càng mang tính toàn cầu (T. Chen et al., 2021). Khối lượng nghiên cứu được công bố đã tăng theo cấp số nhân, với một số lĩnh vực có lượng tài liệu nhân đôi sau mỗi vài năm (Bornmann and Mutz, 2015; Bornmann, Haunschild and Mutz, 2021). Sự bùng nổ thông tin này khiến các cách tiếp cận thủ công truyền thống đối với tổng quan tài liệu và phát triển khái niệm vừa tốn thời gian vừa dễ bỏ sót. Các cách tiếp cận dựa trên AI đã chứng tỏ giá trị đặc biệt trong việc xác định các xu hướng nghiên cứu mới nổi và các khoảng trống nghiên cứu; theo một số nghiên cứu, các tổng quan tài liệu có sự hỗ trợ của AI có thể xác định được nhiều nguồn liên quan hơn so với chỉ dùng các phương pháp truyền thống, và tiết kiệm đáng kể thời gian khi được sử dụng đúng cách (van Dijk et al., 2023). Tuy nhiên, việc giữ một góc nhìn cân bằng và có tính phản biện vẫn là điều then chốt. Dù các công cụ này có thể gợi ý những góc tiếp cận mới hoặc làm lộ ra các khung lý thuyết tiềm năng, chúng chỉ nên được dùng như công cụ hỗ trợ chứ không phải người quyết định hướng nghiên cứu. Một trong những thách thức chính là sự suy giảm tư duy phản biện và tính xác thực (Cardon et al., 2023): suy cho cùng, AI tạo sinh, như tên gọi, có tính ``tạo sinh'' và về mặt thiết kế thì thiên về lặp lại các khuôn mẫu.

Các học giả cũng phải luôn lưu tâm đến các vấn đề về thiên kiến, tính minh bạch và liêm chính trí tuệ (Thorp, 2023). AI tạo sinh phản ánh những thiên kiến xã hội hiện hữu và thường khuếch đại chúng lên nhiều bậc độ lớn (Gorska and Jemielniak, 2023). Việc tích hợp AI tạo sinh vào vòng đời nghiên cứu là một cơ hội đáng kể để nâng cao chất lượng công trình học thuật, với điều kiện được triển khai một cách thận trọng và có đạo đức. Các ứng dụng thực tế bao gồm dùng AI để đề xuất các câu hỏi nghiên cứu ban đầu, xác định các khung lý thuyết liên quan hoặc gợi ý các hướng tiếp cận phương pháp luận. Các nhà nghiên cứu được khuyến khích lưu giữ tài liệu chi tiết về các tương tác với AI, đánh giá một cách phản biện những gợi ý do AI tạo ra, và thường xuyên đối chiếu các phát hiện với các chuẩn mực học thuật đã được thiết lập.

\section{Hình thành ý tưởng và phát triển khái niệm}

Ở giai đoạn sớm nhất của một dự án nghiên cứu, các học giả thường phải ``chăn mèo'': định hướng trong một bức tranh phức tạp gồm những mối quan tâm rộng, kiến thức chưa đầy đủ và các chủ đề đang nổi lên. Quá trình chuyển những ý tưởng rời rạc thành các vấn đề nghiên cứu mạch lạc từ trước đến nay vẫn dựa vào sự sáng tạo, chuyên môn học thuật và việc khảo sát có phương pháp các tài liệu hiện có của nhà nghiên cứu (Creswell, 2014; Amabile, 1983). AI tạo sinh (generative AI) hiện mở ra những con đường tinh vi để tăng cường và đơn giản hóa các quá trình khái niệm ban đầu này, đóng vai trò một công cụ hình thành ý tưởng tiên tiến: làm sáng tỏ các mối liên hệ mới, chỉ ra những vùng chưa được khám phá và giúp kết tinh các khái niệm sơ bộ thành những câu hỏi nghiên cứu tập trung.

Các mô hình ngôn ngữ lớn (LLM) có thể tổng hợp thông tin từ các cơ sở dữ liệu văn bản đồ sộ, qua đó phát hiện những hướng nghiên cứu tiềm năng có thể vẫn bị che khuất trong các tổng quan tài liệu thông thường (Floridi and Chiriatti, 2020; van Dis et al., 2023). Các hệ thống này tạo điều kiện cho điều mà chúng ta có thể gọi là ``hình thành ý tưởng tăng cường'' (augmented ideation), bổ trợ chứ không thay thế đóng góp trí tuệ của nhà nghiên cứu. Các thử nghiệm đương đại cho thấy AI tạo sinh đóng vai trò chất xúc tác cho tư duy sáng tạo: đề xuất các khung lý thuyết thay thế, hé lộ các mối liên hệ liên ngành và thách thức các giả định của nhà nghiên cứu thông qua phân tích phản thực (Rudolph, Tan, and Tan, 2023; Boden, 2009).

Hãy xét một nhà nghiên cứu đang tìm hiểu các hệ quả xã hội của AI trong giáo dục. Thông qua AI tạo sinh, họ có thể khám phá nhiều góc nhìn phân tích đa dạng, từ các nghiên cứu chính sách so sánh giữa các khu vực pháp lý, đến các khảo sát về mô hình đào tạo giáo viên trong việc áp dụng AI, hay các phân tích lịch sử về việc tích hợp công nghệ vào môi trường giáo dục. Hệ thống AI có thể xử lý hiệu quả các tập dữ liệu khổng lồ để nhận diện các xu hướng mới nổi, đặt ra các câu hỏi phương pháp luận và làm nổi bật những phát hiện thực nghiệm gần đây có thể làm giàu khung khái niệm của nhà nghiên cứu. Việc đẩy nhanh giai đoạn khám phá này giúp tiến nhanh hơn tới việc xây dựng giả thuyết. Tuy nhiên, việc vận dụng một cách thận trọng vẫn là điều tối quan trọng. Dù AI tạo sinh xuất sắc trong việc đề xuất các hướng nghiên cứu mới, nó cũng có thể tạo ra những gợi ý lạc đề, thiếu cơ sở lý thuyết hoặc có vấn đề về đạo đức nếu được chấp nhận mà không qua đánh giá phản biện (Thorp, 2023; Lenharo, 2024). Nhà nghiên cứu phải duy trì các chuẩn mực phân tích nghiêm ngặt, coi các đầu ra do AI tạo ra là những nhận định sơ bộ cần được xem xét kỹ lưỡng, kiểm chứng và hoàn thiện. Mục tiêu là tăng cường chứ không phải thuê ngoài quá trình trí tuệ, bảo đảm rằng chuyên môn của con người, kiến thức lĩnh vực và phán đoán học thuật dẫn dắt việc khái niệm hóa nghiên cứu. AI tạo sinh nổi lên như một công cụ cộng tác vô giá trong việc định hình bối cảnh nghiên cứu đương đại. Bằng cách hé lộ những khả năng chưa được khám phá, đặt câu hỏi với các giả định đã được thiết lập và thúc đẩy đối thoại liên ngành, các hệ thống này nâng cao hiệu quả, tính sáng tạo và độ vững chắc của giai đoạn hình thành ý tưởng và phát triển khái niệm. Các phần tiếp theo sẽ xem xét cách những lợi thế ban đầu này bổ trợ cho các quy trình tổng quan tài liệu, phát triển giả thuyết và lập kế hoạch phương pháp luận, từ đó đặt nền tảng cho một cuộc khảo sát thực nghiệm toàn diện.

\section{Khám phá tài liệu và tổng quan nhanh}

Các nhà nghiên cứu đương đại đối mặt với một thách thức chưa từng có trong việc quản lý sự gia tăng theo cấp số nhân của các ấn phẩm học thuật đã nêu ở trên. ``Sự bùng nổ thông tin'' này gây ra những trở ngại đáng kể trong việc xác định các công trình nền tảng, nắm bắt những phát triển tiên tiến nhất và nhận ra các mô thức trí tuệ định hình các lĩnh vực học thuật (Bornmann, Haunschild and Mutz 2021). Các mô hình AI tạo sinh tiên tiến, tận dụng năng lực xử lý ngôn ngữ tự nhiên (NLP) tinh vi, đưa ra những giải pháp mới mẻ để định hướng trong sự phức tạp này, cho phép tiếp cận việc khám phá và tổng quan tài liệu một cách có hệ thống và hiệu quả hơn.

Các công cụ dựa trên AI này đặc biệt xuất sắc trong việc lọc nội dung sơ bộ, xác định hiệu quả các ấn phẩm và bộ dữ liệu quan trọng phù hợp với khung lý thuyết của dự án (Floridi and Chiriatti, 2020). Nhà nghiên cứu có thể tận dụng các mô hình tạo sinh thông qua việc viết prompt (câu lệnh) có mục tiêu để xác định các nghiên cứu nền tảng, nhận diện những học giả có ảnh hưởng và lập bản đồ các cụm khái niệm. Cách tiếp cận này giảm đáng kể thời gian vốn cần cho các tổng quan tài liệu toàn diện, đồng thời có khả năng phát hiện những mối liên hệ giá trị mà phân tích thủ công có thể bỏ sót (van Dis et al., 2023). Hiệu quả của các phương pháp tính toán trong tổng quan hệ thống đã được ghi nhận rõ ràng ở nhiều ngành khác nhau, từ y tế công cộng đến khoa học xã hội (Tsafnat et al., 2014; Baker and Dunbar, 2020).

Một điểm mạnh đặc trưng của các mô hình tạo sinh là khả năng tổng hợp khối lượng lớn công trình học thuật mà vẫn giữ được các sắc thái thiết yếu và các mối quan hệ ngữ cảnh. Tiếp nối ví dụ trước về một nhà nghiên cứu khám phá giao điểm giữa đạo đức AI và giáo dục đại học, thông qua những prompt được soạn kỹ lưỡng, họ có thể nhanh chóng tạo ra các bản tổng quan toàn diện về những cuộc tranh luận chính, từ mối lo ngại về quyền riêng tư đến thiên kiến thuật toán trong tuyển sinh, đồng thời xác định những mảng còn thiếu bằng chứng thực nghiệm hoặc chưa được phát triển về mặt lý thuyết. Cách tiếp cận linh hoạt này là một bước tiến đáng kể so với các phương pháp tổng quan tuyến tính truyền thống, vốn thường đòi hỏi những khoảng thời gian dài để đọc và tổng hợp.

Tuy nhiên, việc tích hợp AI vào tổng quan tài liệu đòi hỏi sự cân nhắc kỹ lưỡng và tính chặt chẽ về phương pháp luận. Các mô hình tạo sinh đôi khi có thể đưa ra thông tin không chính xác hoặc gán sai nguồn, do đó cần kiểm chứng kỹ lưỡng nội dung do AI tạo ra. Mọi tài liệu tham khảo do mô hình AI gợi ý đều phải được kiểm tra và rà soát: không chỉ có lúc chúng không bao hàm đúng nội dung mà mô hình đã nói, mà chúng thậm chí có thể hoàn toàn không tồn tại, dù DOI, khoảng số trang hay đường liên kết trông có vẻ rất hợp lý. Vì vậy, việc duy trì liêm chính học thuật đòi hỏi phải xác thực nguồn một cách hệ thống, kiểm tra cẩn thận các trích dẫn và đánh giá phê phán các bản tổng hợp do AI tạo ra.

AI tạo sinh về cơ bản nâng cao năng lực của nhà nghiên cứu trong việc nhanh chóng xây dựng hiểu biết toàn diện về các lĩnh vực mới hoặc đang phát triển. Bằng cách hợp lý hóa việc xác định và tổng hợp các công trình học thuật liên quan, những công cụ này cho phép phân bổ nguồn lực nghiên cứu một cách chiến lược hơn. Khi được triển khai một cách thận trọng, sự hỗ trợ của AI có thể biến tổng quan tài liệu từ một yêu cầu tốn nhiều thời gian thành một quy trình linh hoạt, lặp lại, tạo ra những hiểu biết sâu sắc hơn và thúc đẩy các hướng nghiên cứu sáng tạo hơn. Tuy nhiên, khó có khả năng trong tương lai gần các tổng quan tài liệu sẽ được thực hiện tự động và không cần giám sát.

\section{Hình thành giả thuyết và các khung khái niệm sơ bộ}

Sự chuyển đổi từ việc khám phá các chủ đề rộng sang những giả thuyết có thể kiểm chứng một cách chính xác là một bước ngoặt quan trọng trong quá trình nghiên cứu. Theo truyền thống, các học giả dựa vào nền tảng lý thuyết, chuyên môn của ngành và quá trình lặp lại có phương pháp để xây dựng giả thuyết (Bryman, 2012; Maxwell, 2013). Hiện nay, AI tạo sinh mang lại những năng lực bổ trợ mạnh mẽ, giúp tăng cường quá trình này mà không thay thế vai trò thiết yếu của nhà nghiên cứu. Các hệ thống tiên tiến này đặc biệt giỏi trong việc gợi ý những mở rộng hợp lý của các lý thuyết đã được xác lập, nhận diện các mối quan hệ biến số tinh tế và làm sáng tỏ các cơ chế nhân quả tiềm năng vốn có thể chưa được khám phá, qua đó thúc đẩy thiết kế nghiên cứu tinh vi hơn.

Hãy xem xét một nhà nghiên cứu đang khảo sát tác động của các chương trình nâng cao năng lực số đối với phát triển cộng đồng. Sau khi phân tích các cấu trúc lý thuyết và kết quả thực nghiệm liên quan, một hệ thống AI tạo sinh có thể đề xuất những giả thuyết tinh tế có tính đến các yếu tố bối cảnh:

Ở những vùng có trình độ học vấn nền thấp hơn, việc đào tạo năng lực số sẽ cho thấy tác động tích cực mạnh hơn đối với hoạt động khởi nghiệp so với những vùng có nguồn lực giáo dục dồi dào.

Gợi ý này buộc các nhà nghiên cứu phải xem xét cách các bối cảnh kinh tế - xã hội điều tiết kết quả của chương trình, có khả năng làm lộ ra những hiệu ứng tương tác phức tạp đáng được nghiên cứu. Tương tự, trong nghiên cứu tổ chức, AI có thể làm nổi bật các cơ chế trung gian cụ thể:

Trong môi trường làm việc từ xa, sự hiện diện của các kênh giao tiếp có cấu trúc đóng vai trò trung gian trong mối quan hệ giữa phong cách lãnh đạo chuyển đổi và sự gắn kết của nhóm.

Điều này thu hút sự chú ý đến những biến số mà ban đầu nhà nghiên cứu có thể đã bỏ sót. Những năng lực này phù hợp với các thực hành nghiên cứu đang phát triển, ngày càng tích hợp phân tích dự báo và các phương pháp dựa trên dữ liệu vào việc hình thành giả thuyết (Shmueli, 2010). Bằng cách phân tích các kho ngữ liệu học thuật đồ sộ, các mô hình ngôn ngữ lớn có thể làm nổi lên những mối liên hệ hợp lý, nêu bật các cấu trúc lý thuyết chưa được khai thác đúng mức và gợi ý các khung phân tích mới. Chẳng hạn, trong nghiên cứu y tế công cộng, một mô hình tạo sinh có thể đề xuất xem xét việc ``niềm tin của cộng đồng đối với các nhà cung cấp dịch vụ chăm sóc sức khỏe'' hoặc ``sự phù hợp văn hóa của thông điệp sức khỏe'' làm trung gian cho hiệu quả của can thiệp ra sao, từ đó khơi gợi các khung khái niệm tinh vi hơn. Sau đó, các nhà nghiên cứu có thể đánh giá một cách có hệ thống các yếu tố này thông qua những phương pháp định tính và định lượng đã được thiết lập, vượt ra ngoài suy đoán lý thuyết ban đầu (Robson and McCartan, 2016; Ciesielska and Jemielniak, 2018). Tuy nhiên, phán đoán học thuật nghiêm ngặt vẫn giữ vai trò tối quan trọng trong quá trình này. Không phải giả thuyết nào do AI tạo ra cũng đứng vững trước sự soi xét về mặt lý thuyết hoặc phù hợp với bằng chứng thực nghiệm hiện có (Maxwell, 2013). Một số gợi ý có thể rơi vào vùng suy đoán hoặc phản ánh những thiên kiến vốn có trong dữ liệu huấn luyện nền tảng. Do đó, các giả thuyết có được nhờ sự hỗ trợ của AI phải được đánh giá cẩn thận dựa trên các khung lý thuyết đã được xác lập, các ràng buộc phương pháp luận và bằng chứng thực nghiệm sẵn có.

Chuyên môn của nhà nghiên cứu, được bồi đắp bởi kiến thức sâu về lĩnh vực và sự nghiêm cẩn học thuật, vẫn là trọng tài tối hậu về tính hợp lệ và ý nghĩa tiềm tàng của một giả thuyết. Việc tích hợp AI tạo sinh một cách chiến lược vào quá trình phát triển giả thuyết có thể nâng cao đáng kể thiết kế nghiên cứu nhờ đẩy nhanh việc xác định các mối quan hệ có thể kiểm chứng, bộc lộ những liên hệ lý thuyết bất ngờ và khuyến khích các khung khái niệm tinh vi hơn. Khi được triển khai đúng cách, việc hình thành giả thuyết với sự hỗ trợ của AI làm phong phú quá trình nghiên cứu, giúp các học giả đi từ những ý tưởng ban đầu đến các mệnh đề có cấu trúc chặt chẽ và có thể kiểm chứng. Cách tiếp cận này kết hợp bề rộng của phân tích do AI dẫn dắt với chiều sâu của chuyên môn con người, cuối cùng củng cố nền tảng cho việc khảo sát thực nghiệm tiếp theo. Kết quả là một thiết kế nghiên cứu vững chắc và tinh tế hơn, được hưởng lợi từ cả năng lực công nghệ lẫn phán đoán học thuật, tạo tiền đề cho những đóng góp khoa học có ý nghĩa.

\section{Hướng dẫn phương pháp luận ban đầu}

Việc lựa chọn các phương pháp nghiên cứu phù hợp, sau khi đã xác lập định hướng nghiên cứu và các giả thuyết sơ bộ, là một điểm quyết định then chốt trong quá trình nghiên cứu. Trước nay, việc chọn phương pháp thường được định hướng bởi các quy ước của ngành, nền tảng lý thuyết và chuyên môn phương pháp luận của nhà nghiên cứu \href{https://paperpile.com/c/uF2NXv/e5GM+47Mj}{(Flick 2018; Gerring 2017)}. Nay, AI tạo sinh có thể đưa ra những chỉ dẫn tinh vi giúp định hướng trong bức tranh phức tạp của các tiếp cận định tính, định lượng và hỗn hợp, qua đó có khả năng cải thiện quá trình ra quyết định về phương pháp luận. Đặc biệt trong các phương pháp hỗn hợp, chẳng hạn phương pháp \emph{Thick Big Data} (Jemielniak, 2020), việc dùng AI để lấp những khoảng trống ở các phương pháp mà ta không phải là chuyên gia có thể thực sự hữu ích và mang lại cách tiếp cận toàn diện hơn đối với hiện tượng được nghiên cứu. Các công cụ AI cũng có thể đề xuất những nghiên cứu thí điểm khoa học dữ liệu mang tính khám phá, vốn có thể được triển khai dễ dàng, một lần nữa nhờ sự hỗ trợ của tác tử lập trình AI.

Hãy xét một nhà nghiên cứu đang tìm hiểu niềm tin của công chúng đối với khoa học và phải lựa chọn giữa phỏng vấn sâu, khảo sát quy mô lớn hoặc một cách tiếp cận tích hợp. Thông qua tương tác với một hệ thống AI tạo sinh, nhà nghiên cứu có thể diễn đạt các yêu cầu then chốt của nghiên cứu, chẳng hạn nhu cầu vừa hiểu bối cảnh vừa có khả năng khái quát hóa thống kê. AI có thể gợi ý một thiết kế phương pháp hỗn hợp tinh vi, bắt đầu bằng các cuộc phỏng vấn định tính để khám phá những câu chuyện về niềm tin trong các cộng đồng khác nhau, sau đó là một khảo sát định lượng để đánh giá mức độ phổ biến của các khuôn mẫu đã nhận diện. Chỉ dẫn này phù hợp với tài liệu phương pháp luận đã được thiết lập, vốn ủng hộ tính bổ trợ của nhiều phương pháp trong việc giải quyết các câu hỏi nghiên cứu phức tạp {[}(Johnson and Onwuegbuzie, 2004; Tashakkori, Teddlie, and SAGE Publications, Inc., 2021){]}.

Như đã đề cập, AI tạo sinh có giá trị trong việc giới thiệu cho nhà nghiên cứu những tiếp cận phương pháp luận chuyên biệt có thể nằm ngoài chuyên môn trực tiếp của họ. Chẳng hạn, một nhà nghiên cứu khảo sát diễn ngôn về chính sách môi trường có thể hưởng lợi từ các gợi ý của AI kết hợp phân tích nội dung định tính truyền thống với các kỹ thuật tính toán tiên tiến như mô hình hóa chủ đề (topic modeling) hoặc phân tích cảm xúc (sentiment analysis). Những khuyến nghị như vậy có thể khuyến khích đổi mới phương pháp luận đồng thời duy trì tính chặt chẽ trong phân tích, kết nối chiều sâu diễn giải với những hiểu biết dựa trên dữ liệu {[}(Flick, 2018){]}.

Đối với các nghiên cứu định tính, AI có thể làm sáng tỏ những đổi mới phương pháp luận mới nổi vượt ra ngoài các tiếp cận thông thường. Nhà nghiên cứu có thể nhận ra khả năng ứng dụng của nghiên cứu hành động có sự tham gia, dân tộc học thị giác hay nhân học số vào bối cảnh nghiên cứu cụ thể của mình. Trong các lĩnh vực định lượng, AI có thể nhận diện những kỹ thuật thống kê hoặc thiết kế thực nghiệm tinh vi phù hợp chính xác với mục tiêu nghiên cứu, chẳng hạn gợi ý các mô hình hiệu ứng hỗn hợp theo chiều dọc cho các nghiên cứu can thiệp giáo dục, hoặc đề xuất các tiếp cận phi tham số vững khi dữ liệu vi phạm các giả định truyền thống.

Tuy nhiên, việc lựa chọn phương pháp luận phải luôn dựa trên các cân nhắc thực tiễn. Dù AI tạo sinh rất giỏi trong việc xác định các tiếp cận phương pháp luận tiềm năng, nhà nghiên cứu vẫn phải đánh giá cẩn trọng các gợi ý đó so với những ràng buộc thực tế, bao gồm chuyên môn sẵn có, nguồn lực, giới hạn về thời gian và các cân nhắc đạo đức. Các khuyến nghị của AI, vốn được rút ra từ các khuôn mẫu trong tài liệu hiện có và dữ liệu huấn luyện, đòi hỏi phải được kiểm chứng cẩn thận trong từng bối cảnh nghiên cứu cụ thể. Các học giả phải vận dụng phán đoán chuyên môn để bảo đảm rằng các phương pháp được chọn phù hợp với mục tiêu nghiên cứu, đồng thời khả thi và đúng đắn về mặt đạo đức trong môi trường nghiên cứu riêng của họ.

Do đó, AI tạo sinh đóng vai trò như một nhà tư vấn phương pháp luận tinh vi trong giai đoạn lập kế hoạch nghiên cứu, làm sáng tỏ các tiếp cận đổi mới đồng thời khuyến khích suy ngẫm phê phán về các lựa chọn phương pháp. Thông qua việc tham gia có chiến lược với các gợi ý do AI tạo ra, nhà nghiên cứu có thể xây dựng những khung phương pháp luận vững chắc, phù hợp bối cảnh và mang tính đổi mới, định hướng hiệu quả cho các giai đoạn tiếp theo của thu thập, phân tích và diễn giải dữ liệu. Sự kết hợp giữa năng lực công nghệ và phán đoán học thuật này làm tăng tiềm năng đổi mới phương pháp luận, đồng thời duy trì những chuẩn mực thiết yếu về tính chặt chẽ học thuật.

\section{Phát triển công cụ tự động và thu thập dữ liệu}

Bước chuyển từ thiết kế nghiên cứu sang thu thập dữ liệu là một thời điểm then chốt, nơi việc triển khai kỹ thuật gặp gỡ kế hoạch phương pháp luận. Trong khi các nhà nghiên cứu theo truyền thống phải đầu tư nhiều thời gian để xây dựng các công cụ nghiên cứu tùy chỉnh, lập trình các công cụ thu thập dữ liệu và quản lý các quy trình tích hợp dữ liệu phức tạp (Munzert et al., 2015), thì AI tạo sinh nay có thể hỗ trợ tinh vi nhằm hợp lý hóa những khía cạnh kỹ thuật này của khâu chuẩn bị nghiên cứu. Những hệ thống tiên tiến này có thể tạo ra mã tùy chỉnh, đề xuất các chiến lược thu thập dữ liệu hiệu quả và hỗ trợ phát triển các quy trình nghiên cứu tích hợp.

Hãy xem xét một nhà nghiên cứu đang tìm hiểu sự tiến triển của chính sách công thông qua các tài liệu của chính phủ. Thay vì tự tay viết các đoạn mã web scraping, họ có thể tận dụng các hệ thống AI tạo mã để xây dựng các quy trình thu thập dữ liệu hiệu quả. Bằng cách cung cấp các tham số cụ thể, chẳng hạn như kho lưu trữ mục tiêu, đặc tả tài liệu và tiêu chí trích xuất, các nhà nghiên cứu có thể nhận được những đoạn mã Python tinh vi để xử lý các yêu cầu web, quản lý phân trang và cấu trúc hóa dữ liệu đầu ra (M. Chen et al., 2021). Việc tự động hóa này đẩy nhanh quá trình phát triển và cho phép nhà nghiên cứu tập trung vào các khía cạnh then chốt như tuân thủ đạo đức, đảm bảo chất lượng dữ liệu và hoàn thiện phương pháp luận.

Khả năng của AI tạo sinh vượt ra ngoài việc thu thập dữ liệu cơ bản, mở rộng sang tiền xử lý tinh vi và chuẩn bị cho phân tích. Chẳng hạn, khi chuẩn bị dữ liệu văn bản cho phân tích tính toán, các nhà nghiên cứu có thể nhận được các đoạn mã do AI tạo ra cho những tác vụ xử lý ngôn ngữ tự nhiên (NLP) nâng cao, bao gồm token hóa, phân tích cú pháp hoặc phân tích cảm xúc. Tương tự, khi kết nối với các API của mạng xã hội, những hệ thống này có thể cung cấp các mẫu mã vững chắc cho xác thực, giới hạn tần suất truy vấn và phân tích cú pháp dữ liệu. Sự hỗ trợ kỹ thuật này giúp các nhà nghiên cứu thu hẹp khoảng cách giữa yêu cầu khái niệm và việc triển khai thực tế một cách hiệu quả hơn.

Khả năng của AI trong việc hỗ trợ tích hợp giữa các môi trường kỹ thuật đa dạng đặc biệt có giá trị. Nghiên cứu hiện đại thường đòi hỏi sự thành thạo nhiều ngôn ngữ lập trình và khung phân tích. AI tạo sinh có thể đề xuất các giải pháp hiệu quả để truyền dữ liệu giữa các nền tảng khác nhau, ví dụ, kết nối các phân tích thống kê dựa trên R với các hệ thống thu thập dữ liệu dựa trên Python. Sự hỗ trợ tích hợp này làm giảm các rào cản kỹ thuật và đẩy nhanh việc phát triển các quy trình nghiên cứu toàn diện.

Do ngày càng nhiều công cụ dựa vào các phương pháp low-code hoặc no-code, việc sử dụng AI cho thu thập dữ liệu và nghiên cứu nhiều khả năng sẽ trở thành một phần tiêu chuẩn trong kho kỹ năng học thuật (Ciechanowski, Jemielniak, and Gloor, 2020).

\section{Kết luận}

AI tạo sinh đã định hình lại sâu sắc quy trình nghiên cứu học thuật, giúp hợp lý hóa việc hình thành ý tưởng, tổng quan tài liệu, xây dựng giả thuyết, lựa chọn phương pháp luận và thu thập dữ liệu. Bằng cách cung cấp các công cụ để khám phá hiệu quả, tinh chỉnh giả thuyết, tư vấn phương pháp luận và tự động hóa kỹ thuật, AI tạo sinh giảm đáng kể thời gian và gánh nặng nhận thức vốn gắn với việc thiết kế nghiên cứu và thu thập dữ liệu. Tuy nhiên, việc tích hợp thành công đòi hỏi sự giám sát quan trọng của con người, bởi những hiểu biết do AI tạo ra luôn phải được kiểm chứng dựa trên tính chặt chẽ lý thuyết, bằng chứng thực nghiệm và các cân nhắc đạo đức. Rốt cuộc, AI tạo sinh nổi lên không phải như một sự thay thế mà như một đối tác có giá trị, tăng cường năng lực của nhà nghiên cứu và giúp họ tập trung sâu hơn vào các khía cạnh sáng tạo và trí tuệ của hoạt động học thuật.
